\documentclass[UTF8,12pt]{ctexart}
\usepackage[a4paper,margin=24mm]{geometry}
\usepackage{amsmath,amssymb,bm,graphicx,adjustbox}
\newcommand{\fitmath}[1]{\begin{adjustbox}{max width=\linewidth}$\displaystyle #1$\end{adjustbox}}
\setlength{\parindent}{0pt}
\setlength{\parskip}{.7em}
\sloppy
\allowdisplaybreaks
\begin{document}
\section*{1998 年考研数学三 · 参考答案}
\subsection*{1998 年真题参考答案}

\subsection*{一、填空题}

（1）\(\displaystyle \dfrac{\displaystyle 1}{\displaystyle e}\)。 （2）\(\displaystyle -\dfrac{\displaystyle \ln x}{\displaystyle x}+C\)，其中 \(\displaystyle C\) 为任意常数。 （3）\(\displaystyle y_t=C(-5)^t+\dfrac{\displaystyle 5}{\displaystyle 12}\left(t-\dfrac{\displaystyle 1}{\displaystyle 6}\right)\)，其中 \(\displaystyle C\) 为任意常数。


（4）\(\displaystyle B=\begin{pmatrix}2&0&0\\0&-4&0\\0&0&2\end{pmatrix}\)。 （5）\(\displaystyle a=\dfrac{\displaystyle 1}{\displaystyle 20},\allowbreak b=\dfrac{\displaystyle 1}{\displaystyle 100}\)；自由度为 \(\displaystyle 2\)。


\subsection*{二、选择题}

（1）D。 （2）B。 （3）C。 （4）B。 （5）A。


三、\(\displaystyle dz=e^{-\arctan(y/x)}[(2x+y)\,dx+(2y-x)\,dy]\)，\(\displaystyle \dfrac{\displaystyle \partial^2z}{\displaystyle \partial x\partial y}=\dfrac{\displaystyle y^2-xy-x^2}{\displaystyle x^2+y^2}e^{-\arctan(y/x)}\)。


四、\(\displaystyle \dfrac{\displaystyle 8}{\displaystyle 15}\)。


五、窖藏 \(\displaystyle t=\dfrac{\displaystyle 1}{\displaystyle 25r^2}\) 年售出，总收入现值最大；当 \(\displaystyle r=0.06\) 时，\(\displaystyle t\approx11\) 年。


六、证明略。（可利用微分中值定理。）


七、（1）\(\displaystyle S_n=\dfrac{\displaystyle 4}{\displaystyle 3}\cdot\dfrac{\displaystyle 1}{\displaystyle n(n+1)\sqrt{n(n+1)}}\)。（2）\(\displaystyle \dfrac{\displaystyle 4}{\displaystyle 3}\)。


八、微分方程为 \(\displaystyle x^2y^{\prime}-3y^2=-2xy\)，则 \(\displaystyle y^{\prime}=3\left(\dfrac{\displaystyle y}{\displaystyle x}\right)^2-\dfrac{\displaystyle 2y}{\displaystyle x}\)。满足给定初始条件的解是 \(\displaystyle y-x=-x^3y\)。


九、（1）\(\displaystyle A^2\) 为 \(\displaystyle n\) 阶零矩阵。（2）矩阵 \(\displaystyle A\) 的特征值全为 0，\(\displaystyle A\) 的属于特征值 0 的特征向量为 \(\displaystyle k_1(-b_2,\allowbreak b_1,\allowbreak 0,\allowbreak \ldots,\allowbreak 0)^{\mathrm T}+k_2(-b_3,\allowbreak 0,\allowbreak b_1,\allowbreak \ldots,\allowbreak 0)^{\mathrm T}+\cdots+k_{n-1}(-b_n,\allowbreak 0,\allowbreak \ldots,\allowbreak 0,\allowbreak b_1)^{\mathrm T}\)，其中 \(\displaystyle k_1,\allowbreak k_2,\allowbreak \ldots,\allowbreak k_{n-1}\) 是不全为零的任意常数。


十、对角矩阵 \(\displaystyle \Lambda=\begin{pmatrix}(k+2)^2&0&0\\0&(k+2)^2&0\\0&0&k^2\end{pmatrix}\)；当 \(\displaystyle k\ne-2\) 且 \(\displaystyle k\ne0\) 时，\(\displaystyle B\) 为正定矩阵。


十一、期望约 14 166.67 元。


十二、（1）\(\displaystyle \dfrac{\displaystyle 29}{\displaystyle 90}\)。（2）\(\displaystyle \dfrac{\displaystyle 20}{\displaystyle 61}\)。

\end{document}
